\section{Tokenized-to-Parametric Transformations}

经验从 Tokenized 载体到 Parametric 载体的转化包含两条路径：Evaluator Internalization 把交互经验固化为评估能力，Policy Internalization 把交互经验固化为决策能力。

\subsection{Evaluator Internalization: Tokenized-to-Evaluator Transformation}

Evaluator Internalization 把 Tokenized agent experience 内化为 Evaluator 参数，使评估能力固化在权重中。产物涵盖 reward model（outcome / process）、verifier、critic 与 failure detector 等形态。纳入前提是 Evaluator 参数被经验数据实质更新；纯 prompt 触发、参数不变的 LLM-as-a-judge 不在此列。

\subsubsection{Outcome-supervised Evaluator Internalization}

Outcome-supervised internalization 把监督信号绑定到一条完整 response、整段轨迹或整个 episode，将其映射为单一评价标签。按输出形态，本节分为判别式与生成式 outcome 评估器两支。

\paragraph{Discriminative Outcome Evaluators}

判别式 outcome 评估器输出标量判定，按信号类型分 binary outcome judgment 与 graded outcome scoring 两类。

Binary outcome judgment 将完整 episode 映射为成功/失败或违规/合规等二值判定。Du et al. \cite{Du23} 把 success detection 写成 visual question answering，用 clip-level 成功标签训练 Flamingo 做二值判定。I-FailSense \cite{Gri25} 扩展到跨任务 failure detection，正例来自原指令、负例来自语义不匹配指令。WebRL \cite{Qi24} 在自演化课程 RL 中训练 outcome reward model，依据指令、历史动作与页面状态判定 web trajectory 是否完成任务。Wen et al. \cite{Wen25b} 把判定目标从 task success 移到 policy violation，对四类规则逐条判违规。

Graded outcome scoring 在二值之外给出更细的信号——分级分值、连续概率或相对偏好。RoboReward \cite{Lee26b} 在真实 rollout 上做 counterfactual relabeling 与 negative clipping，构造 1–5 级 rubric 回归离散 progress。SAFE \cite{Gu25b} 以 VLA policy 的内部 latent 为输入，在 rollout prefix 上学习随时间变化的 failure probability，由 conformal threshold 触发 abort。VLP \cite{Liu25w} 改学相对偏好，从 expert / medium / random 三档轨迹构造比较标签，按 Bradley-Terry 学习 cross-modal preference。

\paragraph{Generative Outcome Evaluators}

生成式 outcome 评估器除 verdict 外还产出自然语言文本，按文本承担的角色分 verdict justification、failure diagnosis 与 corrective critique 三类。

Verdict justification 在判定之外生成支撑该判定的推理。GenRM \cite{Zha24o} 把 reward modeling 写成 next-token prediction，有直出 Yes / No 的 GenRM-Direct 与先生成 CoT rationale 再判的 GenRM-CoT。S2J \cite{Sun25l} 让 judge 判断前先自行解题，把 solving 与 judging reward 联合优化。J4R \cite{Xu25i} 让 judge 在 response order swap 等等价输入上保持一致判断，压低位置偏置。AgentV-RL \cite{Zha26ac} 把 verifier 做成多轮 tool-augmented agent，以 verdict 与 ground truth 的一致性为奖励。CriticGPT-RM \cite{Liu24g} 从不同阶段操作视频自动生成 pairwise analysis 与 preference，下游经 Bradley-Terry 训练独立 reward model。

Failure diagnosis 在判定之外生成对失败的解释或定位。AHA \cite{Dua24} 通过 FailGen 在成功 demonstration 上注入多类失败模式，训练模型同时输出 yes / no 与自由文本 failure rationale。ARMOR \cite{Qi26} 联合训练 detection head 与 reasoning decoder，多轮生成 detection 与 reasoning 后按自信度选取结果。ExeVRM \cite{Son26d} 以执行视频训练，输出 video-level success judgment 与失败定位。

Corrective critique 生成供下游 generator 或 actor 取用的修正性意见。CTRL \cite{Xie25c} 以单测通过为奖励用 GRPO 优化 critique，使 generator 更可能产出通过测试的代码。Co-Evolving Critics \cite{Li26l} 让 critic 给出 diagnosis、actor 据此重写轨迹，critic 奖励含 saturation-aware gain 以防随 policy 进化而失效。RL4F \cite{Aky23} 以 refined output 相对 ground truth 的任务指标为奖励训练 critique generator，是 critique-driven repair 的早期形态。

\subsubsection{Process-supervised Evaluator Internalization}

Process-supervised internalization 把监督信号分配到中间步骤或动作前缀，将局部决策质量写进 Evaluator 参数。按输出形态，本节分为判别式与生成式 process 评估器两支。动因来自一项经验观察：长程 agent 的失败通常是若干局部错误的累积，单一 outcome label 无法把这些错误从一条整体失败的轨迹中分离出来。

\paragraph{Discriminative Process Evaluators}

判别式 process 评估器对每个中间步骤或动作前缀输出标量分。按信号语义分 step correctness scoring、step advantage and progress estimation 与 step-action compatibility 三类。

\textbf{Step correctness scoring} 把每步映射为正确/错误或更细的错误类型标签。Let's Verify Step by Step \cite{Lig23} 在 PRM800K 上用人工 step-level 标签训练 PRM，对每一步预测 positive/negative/neutral。OmegaPRM \cite{Luo24} 用 MCTS 与 binary search 自动定位 first error，以 Monte Carlo 估计替代人工标注，产出 soft process label。PathFinder-PRM \cite{Pal25} 的核心主张是 reward 估计应以错误类型为条件——模型先显式判定该步犯了 math error 还是 consistency error，再基于错误类型预测 step reward，而非从步骤文本直接跳到分数。GUI-Shepherd \cite{Che25h} 把 step correctness 标签落到 GUI 域，人工标注 state-action 对是否正确、GPT-4o 补充 CoT rationale 作为训练增强，推理时输出标量正确性分数。GAIA \cite{Wan26p} 把训练组织为 data flywheel——首轮以正确/错误 action 与 ground truth 比对训练 critic，次轮把 critic 难例回流训练第二轮，持续吸收真实 GUI agent 的错误分布。

\textbf{Step advantage and progress estimation} 在二值之外给出更细的信号：估计一步对未来成功概率的边际贡献或任务进度的连续值。Rewarding Progress \cite{Set24} 从每个 prefix 展开 Monte Carlo rollout 估计未来成功概率的变化量，让 Evaluator 直接学习 step advantage 而非 step correctness。PQM \cite{Li24m} 把问题形式化为 Q-value ranking——用 Plackett-Luce 排序损失替代独立 BCE，强制正确前缀的估计成功概率高于错误前缀。BiRM \cite{Che25s} 在同一 backbone 上联合训练 backward-looking reward head 与 forward-looking value head，同时输出步骤局部正确性与 partial trajectory 仍通向正确答案的概率。AgentPRM \cite{Xi25b} 将这一思路从 reasoning 扩展到多步 agent 决策——用 TD 与 GAE 从 rollout 自动估计 state-action 对上的 Q-value 与 advantage，使 agent trajectory 的中间步骤也能获得 process reward。ProgRM \cite{Zha25z} 在 GUI 轨迹中用 LCS recipe matching 自动定位 key step 并分配归一化 progress label，基于动作历史与当前观测预测任务进度。Agentic-Q \cite{Wan26s} 让终局 success 信号沿 trajectory 反传，训练 step-wise success probability estimator。SWE-Shepherd \cite{Dih26} 在 code 域面临一个特殊困难——代码修改轨迹缺少自然的过程监督信号——于是用相关文件访问、目标文件修改、测试结果变化等启发式指标构造弱 process reward，折扣累积为 step-wise target 训练 repo-level PRM。

\textbf{Step-action compatibility} 在二值正确性之外引入另一种信号：比较同一决策点多个候选动作的相对优劣。V-Droid \cite{Dai25} 的核心机制是 Pairwise Process Preference——不孤立地给单个动作打分，而是强制同一步内正确动作的分数高于若干错误动作，从相对序而非绝对值中学习。IntentScore \cite{Che26s} 把 compatibility 的条件从 state 扩展到 planning intent——action encoder 显式接收 planning intent 作为输入，学习 state-action-intent 三元组的相容性分数，使评估器能区分"动作本身正确"与"动作在当前意图下正确"。

\paragraph{Generative Process Evaluators}

生成式 process 评估器在标量判定之外产出自然语言文本——critique、错误诊断或修正建议。按文本在下游的功能分 step-level critique、verdict with structured reasoning 与 pre-operative corrective feedback 三类。

\textbf{Step-level critique} 对中间步骤生成针对性诊断，指出错误位置、类型与原因，不强制附带数值评分。DeepCritic \cite{Yan25u} 先用 teacher model 生成的 deliberate critique 做 critique teaching SFT，再用 PRM800K 与 MC 自动标注的 step correctness 数据做 critique incentivization GRPO，使模型学会多视角、可纠错的长文本步骤批注。AutoMathCritique \cite{Xi24} 让 actor 在 GSM8K 与 MATH 上生成 flawed reasoning paths，由 GPT-4o 标注首个错误步骤，经 Monte Carlo soft filtering 保留真正有助修正的 critique 做 SFT，产物是可直接驱动 actor refinement 的 step-level 反馈。

\textbf{Verdict with structured reasoning} 在给出判定（正确/错误、有希望/无希望）的同时生成支撑推理，使判定可审计、可质疑。与纯 critique 不同，这类评估器需同时产出 verdict 与 reasoning，且两者之间存在逻辑约束——reasoning 必须支撑 verdict。StepWiser \cite{Xio25c} 让模型在对 reasoning chunk 给出 Positive/Negative verdict 前先写出 meta-reasoning，用 MC continuation 估计的 chunk 前后 Q-value 变化作为监督。OS-Oracle \cite{Wu25m} 通过四类错误注入合成负例，用 CP-GRPO 强化 verdict 与 rationale 之间的一致性。WebArbiter \cite{Zha26ab} 在 reasoning 结构上引入一个关键变化：模型在比较候选动作前必须先归纳适用于当前任务与页面的原则，再从原则推导 preference verdict，而非直接从状态跳到判断。Web-Shepherd \cite{Cha25b} 引入另一种结构——checklist-grounded evaluation——按预定义 checklist 逐项生成 Yes/No/In-Progress 判断，使评估过程被分解为可逐项核验的子问题。SOLE-R1 \cite{Sch26} 把 structured reasoning 扩展到具身时序域，从机器人视频中学习生成 per-timestep progress 分数及解释其判断的 spatiotemporal CoT。PPRM \cite{Xu25h} 在不可验证的多轮 search 任务上按 correctness、relevance、consistency 等原则生成 step score，再通过 ReNorm 与 outcome reward 校准。

\textbf{Pre-operative corrective feedback} 在动作执行前生成诊断与修正建议，使错误在发生前被拦截。GUI-Critic-R1 \cite{Wan25q} 在候选动作执行前依次生成 observation analysis、possible result、critique、correctness score 与 corrective suggestion，用 Suggestion-aware GRPO 强化，使 critic 的建议能实际降低后续执行失败率。

% 所需宏包：booktabs, graphicx
% \usepackage{booktabs}
% \usepackage{graphicx}

\begin{table}[htbp]
\centering
\caption{Overview of Evaluator Internalization methods.}
\label{tab:evaluator-internalization}
\resizebox{\textwidth}{!}{%
\begin{tabular}{@{}lllllll@{}}
\toprule
\textbf{Work} & \textbf{Supervision Granularity} & \textbf{Output Form} & \textbf{Evaluator Output} & \textbf{Label Source} & \textbf{Training Method} & \textbf{Domain} \\
\midrule
% outcome / disc
Du23                            & outcome & disc & binary label              & environment-outcome                    & cls-CE          & embodied        \\
I-FailSense~\cite{Gri25}        & outcome & disc & binary label              & ground-truth match                     & cls-CE          & embodied        \\
WebRL~\cite{Qi24}               & outcome & disc & binary label              & environment-outcome                    & cls-CE          & web             \\
Wen25b                          & outcome & disc & binary label              & rule-based                             & cls-CE          & web             \\
RoboReward~\cite{Lee26b}        & outcome & disc & graded score              & environment-outcome                    & regression      & embodied        \\
SAFE~\cite{Gu25b}               & outcome & disc & graded score              & environment-outcome                    & regression      & embodied        \\
VLP~\cite{Liu25w}               & outcome & disc & preference                & environment-outcome                    & ranking         & embodied        \\
\addlinespace[3pt]
% outcome / gen
GenRM~\cite{Zha24o}             & outcome & gen  & structured verdict        & ground-truth match                     & gen-CE          & text/reasoning  \\
S2J~\cite{Sun25l}               & outcome & gen  & structured verdict        & ground-truth match                     & RL-PG           & text/reasoning  \\
J4R~\cite{Xu25i}                & outcome & gen  & structured verdict        & ground-truth match                     & RL-PG           & text/reasoning  \\
AgentV-RL~\cite{Zha26ac}        & outcome & gen  & structured verdict        & ground-truth match                     & gen-CE + RL-PG  & text/reasoning  \\
CriticGPT-RM~\cite{Liu24g}      & outcome & gen  & preference                & environment-outcome$^\dagger$          & gen-CE          & embodied        \\
AHA~\cite{Dua24}                & outcome & gen  & failure diagnosis         & rule-based                             & gen-CE          & embodied        \\
ARMOR~\cite{Qi26}               & outcome & gen  & failure diagnosis         & environment-outcome + human            & cls-CE + gen-CE & embodied        \\
ExeVRM~\cite{Son26d}            & outcome & gen  & failure diagnosis         & environment-outcome                    & gen-CE          & GUI             \\
CTRL~\cite{Xie25c}              & outcome & gen  & corrective suggestion     & environment-outcome                    & RL-PG           & code            \\
ECHO~\cite{Li26l}               & outcome & gen  & corrective suggestion     & environment-outcome                    & RL-PG           & text/reasoning  \\
RL4F~\cite{Aky23}               & outcome & gen  & corrective suggestion     & ground-truth match                     & RL-PG           & text/reasoning  \\
\midrule
% process / disc
Let's Verify~\cite{Lig23}       & process & disc & binary label              & human                                  & cls-CE          & text/reasoning  \\
OmegaPRM~\cite{Luo24}           & process & disc & binary label (soft)       & MC-rollout                             & cls-CE          & text/reasoning  \\
PathFinder-PRM~\cite{Pal25}     & process & disc & binary label (error-typed)& human + MC-rollout                     & cls-CE          & text/reasoning  \\
GUI-Shepherd~\cite{Che25h}      & process & disc & graded score              & human + teacher-model                  & cls-CE          & GUI             \\
GAIA~\cite{Wan26p}              & process & disc & binary label              & ground-truth match                     & cls-CE          & GUI             \\
Rewarding Progress~\cite{Set24} & process & disc & progress/advantage        & MC-rollout                             & regression      & text/reasoning  \\
PQM~\cite{Li24m}                & process & disc & progress/advantage        & MC-rollout                             & ranking         & text/reasoning  \\
BiRM~\cite{Che25s}              & process & disc & progress/advantage        & MC-rollout                             & cls-CE + regression & text/reasoning \\
AgentPRM~\cite{Xi25b}           & process & disc & progress/advantage        & MC-rollout (TD/GAE)                    & regression      & web             \\
ProgRM~\cite{Zha25z}            & process & disc & progress/advantage        & ground-truth match (LCS)               & regression      & GUI             \\
Agentic-Q~\cite{Wan26s}         & process & disc & progress/advantage        & environment-outcome                    & regression      & GUI             \\
SWE-Shepherd~\cite{Dih26}       & process & disc & progress/advantage        & rule-based                             & regression      & code            \\
V-Droid~\cite{Dai25}            & process & disc & preference                & human                                  & ranking         & mobile          \\
IntentScore~\cite{Che26s}       & process & disc & graded score              & ground-truth match                     & ranking         & GUI             \\
\addlinespace[3pt]
% process / gen
DeepCritic~\cite{Yan25u}        & process & gen  & step critique             & teacher-model + MC-rollout             & gen-CE + RL-PG  & text/reasoning  \\
AutoMathCritique~\cite{Xi24}    & process & gen  & step critique             & teacher-model + MC-rollout             & gen-CE          & text/reasoning  \\
StepWiser~\cite{Xio25c}         & process & gen  & structured verdict        & MC-rollout                             & RL-PG           & text/reasoning  \\
OS-Oracle~\cite{Wu25m}          & process & gen  & structured verdict        & rule-based (error injection)           & RL-PG           & GUI             \\
WebArbiter~\cite{Zha26ab}       & process & gen  & structured verdict        & ground-truth match + teacher-model     & gen-CE + RL-PG  & web             \\
Web-Shepherd~\cite{Cha25b}      & process & gen  & structured verdict        & human                                  & gen-CE          & web             \\
SOLE-R1~\cite{Sch26}            & process & gen  & structured verdict        & environment-outcome                    & gen-CE + RL-PG  & embodied        \\
PPRM~\cite{Xu25h}               & process & gen  & structured verdict        & environment-outcome                    & gen-CE          & text/reasoning  \\
GUI-Critic-R1~\cite{Wan25q}     & process & gen  & corrective suggestion     & ground-truth match                     & gen-CE + RL-PG  & GUI             \\
\bottomrule
\end{tabular}%
}
\end{table}

\subsubsection{Discussion}

Outcome label 易于从环境信号、执行结果或人工裁判得到，对具有明确成功标准的任务可直接对齐最终目标，在具身的 episode-level success detection 上几乎是默认选择，下游取用方式直接落在 scalar reward。Process 判别式围绕 credit assignment 快速扩张，演化轴是从 correctness 到 advantage、progress、Q-value、bidirectional reward 的语义细分，再向 GUI、computer-use、code、embodied 等带显式中间状态的领域落地；代价是 step-level label 的获取成本（要么人工标注，要么 MC 自动估计）与噪声。生成式两支让 Evaluator 输出可读 critique 或 diagnosis，下游可塌缩为标量进入 Evaluator-Driven Optimization 的 RL、可取 binary 做 Best-of-N 筛选、也可保留全文驱动 step-wise refinement。生成式评估模型的代价是 critique 的正确性缺乏自动验证，比如一条听起来合理的 critique 可能与环境真实反馈不一致。

跨监督粒度看，process 监督换更细的 credit assignment，付出标注成本与噪声；跨产物形态看，生成式以更高推理成本换更丰富的下游 utilization。两条轴上的选择由下游任务的 verifiability 与监督预算共同决定，不存在普适最优。

\subsection{Policy Internalization: Tokenized-to-Policy Transformation}

Policy Internalization 把 Agent Trajectory，比如自然语言轨迹、工具调用记录、代码编辑历史、GUI 操作序列、reasoning chains， 内化为 Policy 参数，使序贯决策能力固化在权重中。内化后，模型在后续任务中无需检索、拼接或重放原始经验即可决策。按写入 Policy 参数的方式分为三类：imitation-based、environment-reward-based 与 preference-based。

本节的纳入边界由进入 policy update 的监督信号来源决定。信号为非参数化的环境时，比如环境可验证结果、unit test 通过、ground-truth match、rule-based verifier等，归本节；信号来自参数化 evaluator（trained RM / PRM / LLM-as-a-judge）时，不归本节。

\subsubsection{Imitation-based Policy Internalization}

Imitation-based internalization 把高质量 agent trajectories 当作监督示范，通过 behavior cloning、SFT 或 rejection-sampled fine-tuning 把示范动作序列写入 Policy 参数。

主线是 success-filtered behavior cloning。\cite{Pat24} 让 web agent 在 WebArena 中自生成 action-observation trajectories，以环境报错与 self-critique 滤出可信动作，在 QLoRA 下做 autoregressive SFT；SWE-Gym \cite{Pan24} 把过滤信号换成可执行 unit test，以测试通过判定 patch trajectory 成败后做 rejection-sampling fine-tuning。命名有时偏离 objective：AppVLM \cite{Pap25b} 称 RFT / ReST，但最终更新是 success-only MLE，失败轨 return = 0 且不进训练集；具身 VLA 的 \cite{Guo25d} 称 online RL，但 RL 阶段只在 action head 上做探索，最终全模型更新是对 expert 与新发现成功轨并集的 supervised learning。当过滤信号来自参数化 evaluator、但 evaluator 只选数据不进 loss 时，工作仍属 imitation：Go-Browse \cite{Gan25b} 的成功标签由 VLM-as-a-judge 给出、仅对成功轨 SFT，NNetNav \cite{Mur24b} 经 hindsight relabeling 改写交互后用 ORM 过滤轨迹，AndroidGen \cite{Lai25d} 用 StepCritic 切出成功子段再做 LoRA-SFT。

示范来源可来自更强的 teacher。Explorer \cite{Pah25} 用 GPT-4o 多 agent pipeline 合成 multimodal web trajectories，经 Task Verifier 过滤后让学生模型学 next action；AgentTuning \cite{Zen23d} 把多环境 teacher trajectories 与通用 instruction data 混合，使模型在习得 agent 行为协议的同时保持 general instruction-following。

\subsubsection{Environment-Reward-based Policy Internalization}

Environment-Reward-based internalization 让 agent 与环境交互产生 trajectories，用非参数化、可程序验证的反馈直接进入 policy optimization 的 objective。可验证信号涵盖任务成败、unit test pass / fail、execution output、SQL execution correctness、ground-truth answer match、GUI 任务完成状态、机器人操作成败等环境本身可判定的结果；成功与失败轨迹都进入优化，reward 把前者转为正信号、后者转为负信号。困难出现在长程多轮的真实环境：一条数十步的轨迹只在终局拿到单个稀疏 verifiable reward，信号难以利用。相关工作围绕这一困难展开。

% 基本可行性由两项早期工作确立。GLAM \cite{Car23} 把 FLAN-T5 放入 BabyAI-Text、以 sparse goal reward 经 PPO 在线 grounding；LLaRP \cite{Szo23} 连接冻结的 LLaMA 与视觉编码器、以 success / subgoal / invalid-action penalty 组成的环境 reward 经 DD-PPO 训练 adapter 与 heads。动作可执行、环境能返回可验证结果时，Tokenized trajectory 即可经 RL 直接转化为 Policy 参数。

\textbf{Exploration preservation.} 长程 agent 在训练早期极易坍缩到低探索策略，一旦坍缩便不再产生成功轨迹、梯度随之消失。EPO \cite{Xu25k} 在 PPO 或 GRPO 外层加 trajectory-level entropy regularization 与 entropy smoothing，缓解 sparse-reward multi-turn RL 的 exploration-exploitation cascade failure。沿熵调控扩展的有：Agentic Reinforced Policy Optimization \cite{Don25d} 在高不确定性 tool step 触发局部 branching，AEPO \cite{Don25c} 把 entropy 同时纳入 advantage 与 clipping 并动态分配 global 与 branch rollout budget，EMPG \cite{Wan25ad} 用当前 entropy 缩放每步 advantage、再加鼓励进入更低熵后继状态的 future-clarity bonus。另一条思路是保证成功信号不断流：ARPO \cite{Lu25f} 在 GUI 环境维护成功 trajectory replay buffer，采样组全失败、无梯度时从缓冲注回正样本；SoLS \cite{Pap25c} 在 mobile app control 中对正优势样本强更新、对负优势样本做 PPO-like 保守正则，并配合成功转移回放 STR；AgentGym-RL \cite{Xi25c} 的 ScalingInter-RL 用先短后长的 horizon curriculum，让 agent 先在可完成的短程任务上拿到信号。

\textbf{Fine-grained credit assignment.} 归因粒度从整轨持续下沉到 turn、step、token。WebAgent-R1 \cite{Wei25} 以 M-GRPO 并行采样 browser trajectories、把 binary success reward 转成 group-relative advantage；\cite{Gol25} 在 SWE 环境以 unit test pass / fail 加 trajectory length penalty 做 token-level DAPO。actor-critic 一支用 learned value 估计中间 advantage：ArCHer \cite{Zho24f} 先用环境 reward 训练 utterance-level Q critic、再把 multi-turn advantage 传给 token actor，Turn-PPO \cite{Li25ae} 把多轮交互建模为 turn-level MDP、在 turn 粒度用 critic 做 GAE。critic-free 一支从 group 比较中估计相对优势：GiGPO \cite{Fen25c} 用 anchor-state grouping 比较不同 rollout 在相同中间状态下的动作、同时构造 episode-level 与 step-level relative advantage，LOOP \cite{Che25af} 在 AppWorld API 环境用 leave-one-out PPO 做 token-level 更新。另有把信号挂到中间语义的做法：IGPO \cite{Wan25y} 用每轮交互前后 ground-truth answer log-prob 的增量定义 information-gain intrinsic reward，把 observation 的真实信息贡献转成 turn-level dense 信号；RLVMR \cite{Zha25ao} 给轨迹挂 \texttt{<planning>}、\texttt{<explore>}、\texttt{<reflection>}、\texttt{<monitor>} 标签，用 rule-verifiable process reward 加 outcome reward 让 meta-reasoning 行为可被显式监督。

\textbf{Verifiable reward densification.} 改造 reward 信号本身：从环境挖掘可程序验证的中间结构，把单一终局 flag 写成 dense composite reward，使密集信号在每步可得。ToolRL \cite{Qia25b} 把 tool-use reward 拆成 format correctness 与 tool name / arg name / arg value 三层匹配。text-to-SQL 域最集中：SQL-Trail \cite{Hua26d} 在 reason-execute-observe 循环里用 execution、turn-budget、schema grounding 等六项组成 composite reward，SQL-ASTRA \cite{Li26r} 的 Column-Set Matching Reward 从结果表列值集合比较给 dense partial-credit、再用 Asymmetric Trajectory Reward 聚合整轨。LongNav-R1 \cite{Hu26e} 在长程 navigation 中用 geodesic progress shaped reward 加 terminal SPL、以 critic-free time-aware baseline 做 REINFORCE++ 风格更新。GUI 与 VLA 域把 rule-based verifiable reward 直接送入 GRPO：GUI-R1 \cite{Luo25b}、AgentCPM-GUI \cite{Zha25an}、SimpleVLA-RL \cite{Li25aa} 在 format 合法性之上叠加 action-type、click-point、text-match、bbox 命中等可判定项，UI-S1 \cite{Lu25j} 在 semi-online RL 中用模型自生成历史、偏离 expert 时由 patch module 续接，reward 由 offline ground-truth matching 提供。VLM as Decision-Making Agents \cite{Zha24s} 把 CoT 加 action 当作端到端 text rollout，用 λ 重加权 reasoning token 的 log-prob 贡献后以环境 reward 做 PPO。同一范式被扩到多任务与泛化研究：AgentRL \cite{Zha25ag} 以 deterministic verifiable rewards 加 task advantage normalization 稳住多任务异步 RL，Generalization in Mobile Agents \cite{Gu26b} 在异步 emulator 上验证 instance、template、app 三层泛化。Q-SFT \cite{Hon24} 处在本子类的形式边界——形式接近 SFT，但 sample weight 由 Bellman target / Q-value 给出，机制更接近 offline value-based RL，reward provenance 仍是环境信号。

\subsubsection{Preference-based Policy Internalization}

Preference-based internalization 从 agent experience 中构造轨迹间或步骤间的相对偏好，用 DPO 或其它 trajectory-pair optimization 直接更新 Policy。源信号是相对序：优势项来自成功任务、通过验证的执行、修正后的轨迹或更高效的操作路径，劣势项来自失败尝试、错误工具调用或未通过测试的代码修改。偏好的粒度从整轨下沉到轨迹内部的局部单元。

\textbf{Trajectory-level preference.} 偏好定义在整条轨迹上。DMPO \cite{Shi24c} 把 DPO 从单轮 response preference 扩到 multi-turn trajectory preference，引入 state-action occupancy measure 与 length normalization 处理长轨迹下的 preference objective。ETO \cite{Son24} 先用 expert trajectory 做 behavior cloning、再让 policy 主动探索，把成功示范与探索失败轨配成 failure-success pair 持续 DPO。Agent-RLVR \cite{Da25} 在 SWE 环境用 guidance 帮助模型探索更可能成功的修复轨迹，再依 unit test pass / fail 构造 preference pairs 做 DPO——guidance 在此是数据生成器，不进 reward。

\textbf{Sub-trajectory preference.} 偏好下沉到轨迹内部的局部单元。\cite{Che24k} 在 tool-use inference tree 中把"同一父节点下成功 step 优于错误分支 step"作为 step-level preference pair，在 SFT 后接续 DPO。HPL \cite{Gao25c} 同时引入 trajectory、step、action group 三层偏好，group 由 fixed、entropy 或 semantic segmentation 划出，最终目标是 BC 与三层 DPO 的复合，其中 group-DPO 是主要增益来源。WEPO \cite{Liu25z} 把偏好收缩到 web element 层面，用目标元素与 DOM 距离近的干扰元素做对比训练 element selection policy，标签来自 benchmark GT 与 DOM 结构，不依赖 learned judge。

% 所需宏包：booktabs, graphicx
% \usepackage{booktabs}
% \usepackage{graphicx}

\begin{table}[htbp]
\centering
\caption{Overview of Policy Internalization methods.}
\label{tab:policy-internalization}
\resizebox{\textwidth}{!}{%
\begin{tabular}{@{}lllllll@{}}
\toprule
\textbf{Work} & \textbf{Learning Objective} & \textbf{Credit Granularity} & \textbf{Supervision Signal} & \textbf{Optimization Algorithm} & \textbf{Experience Source} & \textbf{Domain} \\
\midrule
% imitation
\cite{Pat24}   & imitation  & trajectory & model-judge        & SFT/BC                              & self          & web                                   \\
\cite{Lai25d}  & imitation  & step       & model-judge        & SFT/BC (LoRA)                       & teacher       & mobile                                \\
\cite{Pap25b}  & imitation  & trajectory & task-outcome       & SFT/BC                              & human / self  & mobile                                \\
\cite{Pah25}   & imitation  & trajectory & model-judge        & SFT/BC                              & teacher       & web                                   \\
\cite{Gan25b}  & imitation  & trajectory & model-judge        & SFT/BC                              & self          & web                                   \\
\cite{Zen23d}  & imitation  & token      & task-outcome       & SFT/BC                              & teacher       & tool-use                              \\
\cite{Mur24b}  & imitation  & trajectory & model-judge        & SFT/BC                              & self          & web                                   \\
\cite{Pan24}   & imitation  & trajectory & unit-test          & rejection-sampling                  & teacher       & code                                  \\
\cite{Guo25d}  & imitation  & trajectory & task-outcome       & SFT/BC                              & human / self  & embodied                              \\
\addlinespace[3pt]
% preference
\cite{Shi24c}  & preference & trajectory & ground-truth match & DPO                                 & teacher       & web / embodied / text-reasoning       \\
\cite{Son24}   & preference & trajectory & task-outcome       & DPO                                 & teacher / self& web / embodied / text-reasoning       \\
\cite{Da25}    & preference & trajectory & unit-test          & DPO                                 & self          & code                                  \\
\cite{Che24k}  & preference & step       & model-judge        & DPO (SFT warm-start)                & teacher       & tool-use                              \\
\cite{Gao25c}  & preference & step       & task-outcome       & DPO (traj/step/group) + BC          & teacher       & embodied / web / SQL                  \\
\cite{Liu25z}  & preference & step       & ground-truth match & DPO                                 & human         & web                                   \\
\midrule
% env-reward
\cite{Car23}   & env-reward & step       & task-outcome       & PPO                                 & self          & embodied                              \\
\cite{Szo23}   & env-reward & step       & shaped-composite   & DD-PPO                              & self          & embodied                              \\
\cite{Zho24f}  & env-reward & token      & task-outcome       & hierarchical actor-critic (ArCHer)  & self          & text-reasoning / web                  \\
\cite{Wei25}   & env-reward & trajectory & task-outcome       & M-GRPO                              & self          & web                                   \\
\cite{Gol25}   & env-reward & trajectory & unit-test          & DAPO                                & self          & code                                  \\
\cite{Luo25b}  & env-reward & step       & ground-truth match & GRPO                                & self          & GUI                                   \\
\cite{Zha25an} & env-reward & step       & ground-truth match & GRPO                                & self          & mobile                                \\
\cite{Li25aa}  & env-reward & trajectory & task-outcome       & GRPO                                & self          & embodied                              \\
\cite{Zha24s}  & env-reward & token      & task-outcome       & PPO                                 & self          & text-reasoning / embodied             \\
\cite{Xi25c}   & env-reward & trajectory & task-outcome       & GRPO (ScalingInter-RL)              & self          & web / embodied / tool-use             \\
\cite{Don25d}  & env-reward & token      & shaped-composite   & ARPO                                & self          & tool-use                              \\
\cite{Zha25ag} & env-reward & token      & rule-verifiable    & GRPO                                & self          & embodied / web / tool-use             \\
\cite{Wan25y}  & env-reward & turn       & ground-truth match & GRPO                                & self          & web                                   \\
\cite{Lu25j}   & env-reward & step       & shaped-composite   & semi-online RL                      & human         & GUI                                   \\
\cite{Xu25k}   & env-reward & trajectory & task-outcome       & EPO (on PPO/GRPO)                   & self          & embodied / text-reasoning             \\
\cite{Hua26d}  & env-reward & trajectory & shaped-composite   & GRPO                                & self          & SQL                                   \\
\cite{Hu26e}   & env-reward & turn       & task-outcome       & HAPO                                & self          & embodied                              \\
\cite{Don25c}  & env-reward & token      & shaped-composite   & AEPO                                & self          & tool-use                              \\
\cite{Wan25ad} & env-reward & step       & shaped-composite   & EMPG (on GRPO/DAPO)                 & self          & web / embodied / tool-use             \\
\cite{Lu25f}   & env-reward & trajectory & rule-verifiable    & GRPO                                & self          & GUI                                   \\
\cite{Fen25c}  & env-reward & step       & task-outcome       & GiGPO                               & self          & embodied / web / tool-use             \\
\cite{Qia25b}  & env-reward & step       & shaped-composite   & GRPO                                & self          & tool-use                              \\
\cite{Gu26b}   & env-reward & trajectory & task-outcome       & GRPO                                & self          & mobile                                \\
\cite{Che25af} & env-reward & token      & task-outcome       & LOOP                                & self          & tool-use                              \\
\cite{Li25ae}  & env-reward & turn       & task-outcome       & Turn-PPO                            & self          & web / text-reasoning                  \\
\cite{Hon24}   & env-reward & token      & task-outcome       & Q-SFT                               & human$^\dagger$ & text-reasoning / web / embodied     \\
\cite{Li26r}   & env-reward & step       & shaped-composite   & GRPO                                & self          & SQL                                   \\
\cite{Zha25ao} & env-reward & step       & rule-verifiable    & GRPO                                & self          & embodied / text-reasoning             \\
\cite{Pap25c}  & env-reward & step       & task-outcome       & SoLS (+STR)                         & self          & mobile                                \\
\bottomrule
\end{tabular}%
}
\end{table}


\subsubsection{Discussion}

三类的差别落在 policy update 如何利用一条轨迹：只取其中的成功动作、取带正负号的全部 reward、还是取轨迹之间的相对序。这一选择同时设定各自的信号承载力、训练稳定性与能力上限，三者并非独立旋钮。Imitation-based 以成功动作序列做最大似然，cold-start 阶段能最快教会模型动作空间、输出协议与基本交互流程；正因监督只来自成功示范，示范覆盖之外的策略就无从经 MLE 习得，有限成功路径会把模型留在表层格式或局部启发式，失败轨被整体弃用。Environment-Reward-based 让 reward 同时携带成功与失败的符号，于是试错能发现示范之外的新策略、能力上限最高；同一个开放试错的设定也意味着 reward 稀疏时 credit assignment 困难、online rollout 昂贵、reward hacking 难防，且 RL 的高方差梯度一旦坍缩或被 hack，所损伤的能力不像编辑一条文本规则那样可逐条回退。Preference-based 取轨迹间的相对序：相对序是 supervised 目标，无需估计高方差 policy gradient，训练比 RL 稳，也把失败轨作为 negative 纳入而非弃用；代价是相对序只约束成对轨迹的高下、不提供绝对信号，效果因此压在 preference pair 的构造质量上，trajectory-level 偏好仍可能缺乏细粒度归因。

三类常被接成同一条 pipeline 的不同阶段：imitation 立起格式与协议，reward optimization 或 preference learning 接手能力突破——它们之间因此更多是互补而非互斥（此类阶段串接的 integration 分析见 §7）。Environment-Reward 与 Preference 沿同一条 credit-assignment 轴下沉——前者向 turn / step / token，后者向 step / element / group——分野在监督形式：绝对 verifiable reward 进 RL objective，相对 pair 进 supervised loss。三类还共享一个 Parametric 载体的属性：经验写进权重后不可局部寻址，新增或修正一类行为需要重新训练，而非像 Tokenized 载体那样逐条编辑。
